\section{Limitations and Threats to Validity}
\label{sec:limitations}

\Cref{tab:limitations} lists the remaining limitations of this study, how each could affect the conclusions, and
how it is addressed. None of them changes the qualitative conclusions: the gap between the DDPM and the VAEs
exceeds every source of uncertainty by more than an order of magnitude in FID.

\begin{table}[H]
\centering\small
\caption{Remaining limitations, their effect, and how they are addressed.}
\label{tab:limitations}
\begin{tabularx}{\textwidth}{@{}>{\raggedright\arraybackslash}p{3.3cm}>{\raggedright\arraybackslash}X>{\raggedright\arraybackslash}X@{}}
\toprule
Limitation & Effect on the conclusions & How it is addressed\\
\midrule
\textbf{Shortened DDPM training} (150k of 800k steps) &
Our DDPM numbers are a \emph{lower bound} on what the implementation reaches; FID was still improving
(\cref{fig:fid-train}). It can only understate the DDPM's lead. &
Disclosed as the only deviation from the paper (\cref{tab:hparams}). The likelihood bound already matches the paper
to within $0.04$\,bits/dim. Checkpoints are resumable, so training can be continued.\\
\addlinespace
\textbf{FID with 10k samples} (papers: 50k) &
FID is biased upward at smaller $N$, so our values are not directly comparable to published FID-50k. &
The protocol is stated explicitly (\cref{sec:eval}) and applied identically to every model. KID, which is unbiased
at this $N$, gives the same ranking between VAE and DDPM. Evaluation noise is measured on three independent sample
sets.\\
\addlinespace
\textbf{Final checkpoint, not best} &
The paper reports its best FID over training, which is slightly optimistic. &
Our choice was fixed in advance and avoids selection bias.\\
\addlinespace
\textbf{One training run per model} &
The $\pm$ values cover evaluation noise, not variance across training seeds. Small differences (e.g.\ between the
two VAEs) should be read cautiously. &
Conclusions rely only on differences far larger than the measured noise.\\
\addlinespace
\textbf{Classifier-based thematic metrics} &
They depend on one classifier, which behaves unpredictably on images far from real data (\cref{tab:thematic}). &
Interpreted cautiously, and cross-checked by the blind rating.\\
\addlinespace
\textbf{Automated blind rater} &
A single vision-language model is not a human study, and the answer key was accessible to it. &
Reported as supporting evidence only. The answers contradict the key where a blind rater errs, and a human rating
tool is included in the repository.\\
\addlinespace
\textbf{Original VAE method} &
Stronger VAE variants (hierarchical VAEs, learned priors) would narrow the gap. &
Out of scope by design: the task asks for the original approaches.\\
\addlinespace
\textbf{Training not compiled} &
Our benchmark measured a $1.5\times$ faster DDPM training step with \texttt{torch.compile} ($16.1$ vs.\
$10.6$\,steps/s), which was not enabled for the reported runs. This affects the training-time figures only, not the
models or results. &
Reported honestly as measured. Compilation is one configuration flag (\texttt{train.compile}) for future runs.\\
\addlinespace
\textbf{Code availability} &
A results report is only verifiable if the code is accessible. &
All code, configurations, logs, raw evaluation outputs and a reproduction guide are public at \repo.\\
\bottomrule
\end{tabularx}
\end{table}
